\section{Methodology}\label{sec:6}

The evaluation is designed to ask, for each design principle of \secref{sec:3}, what measurable property operationalizes it and what the system reports for that property.
We describe the corpus, the recommendation baselines, the explanation-generation procedure, the calibration procedure, the counterfactual probe, and the adversarial probe; we map each to a user-facing property; we close with the user-study stance and its justification.
Numerical results are reported in \secref{sec:7}.

\leadpara{Corpus and labeling}
The evaluation uses a fixed demo corpus of 30 resumes, 15 job descriptions, and 47 manually labeled resume--job relevance pairs, with an additional 21 strong and 26 partial labels used for the calibration set.
The corpus is frozen and committed; no evaluation run mutates it.
The relevance scale is a four-point ordinal scale (irrelevant, partially relevant, relevant, strongly relevant), annotated by two independent reviewers with reconciliation on disagreement.
Macro-averaged metrics are reported across all 30 resumes so that no single profile dominates the score.
The corpus is small by the standards of production recommendation systems and is reported as such; it is large enough to support the controlled, end-to-end probes used here, and it is what the prototype is currently fit to evaluate.
We return to the implications of corpus size in \secref{sec:8}.

\leadpara{Recommendation quality (operationalizes G1 and G4)}
Recommendation quality is reported as \ndcg{} over the labeled pairs at cutoff $K{=}5$~\cite{jarvelin2002ndcg}.
We compare lexical baselines (BM25, TF--IDF)~\cite{robertson1995okapi}, dense semantic similarity using a sentence-transformer encoder (all-MiniLM-L6-v2, $d{=}384$)~\cite{reimers2019sentence}, skill overlap, hybrid semantic--skill rankers, fixed multimodal blending with the same six channels the user sees in the explanation panel, learned logistic-regression fusion, and reciprocal rank fusion across four list views~\cite{cormack2010reciprocal}.
The user-facing interpretation of this metric is the relevance of the top five ranked items on average, which is the property a user reading a list experiences.
We additionally report the precision and recall at $K{=}5$ to surface the trade-off between exhaustiveness and accuracy.

\leadpara{Explanation faithfulness (operationalizes G2)}
Explanation faithfulness is measured by the proportion of explanation bullets whose named component is consistent with the score that the named component contributes to the ranking.
A bullet that names \emph{skill overlap} but is generated when the skill-overlap component is below a configured threshold is counted as inconsistent.
The rule-based explainer in the prototype is the primary generator; an LLM-based explainer is also evaluated for comparison, and its outputs are reported in the supplementary material.
The metric is reported together with explanation specificity (the proportion of bullets that name a concrete field rather than a general statement) and explanation consistency (the proportion of synthetic profile variants under which the explanation does not change without the underlying input changing).
The user-facing interpretation of these metrics is the user's ability to act on an explanation: a faithful, specific, and consistent explanation is one a user can trust to point at a real change they might make.

\leadpara{Confidence calibration (operationalizes G3)}
Confidence calibration is reported as the \ece{} between the system's reported confidence and the empirical rate at which the top-ranked item is among the labeled relevant items, computed in ten equal-width confidence bins.
The uncalibrated system is a composite of the same six channels that drive the explanation; the calibrated system applies Platt scaling to the composite score~\cite{platt1999probabilistic,guo2017calibration} and reports a confidence value alongside the ranked list.
We additionally report the Brier score as a complementary measure~\cite{brier1950verification} and the calibration parameters (slope and intercept) of the Platt scaler.
The user-facing interpretation of \ece{} is the user's ability to know when the system is sure and when it is not: a well-calibrated system lets a user treat a 0.9 confidence as a strong signal and a 0.5 confidence as a weak one, with predictable consequences.

\leadpara{Counterfactual probe (operationalizes G3)}
The counterfactual probe generates, for each labeled pair, a small set of single-field edits that would change the ranking, and reports the proportion of edits that, when applied, change the rank by at least one position.
The probe is a proxy for the counterfactual view of \figref{fig:portal}(d); it answers the question of whether the system's reported reasons are sufficient to predict the effect of a small edit.
The probe is restricted to edits within a pre-defined edit space (skills the user has but did not list, experience tiers the user could plausibly attain within the evaluation window, and remote-policy toggles) to avoid counterfactuals that have no actionable interpretation.
The user-facing interpretation is the user's ability to use an explanation as a guide to action: if the explanation names skill overlap as the reason for a low rank, a single-skill addition should move the rank.

\leadpara{Adversarial probe (operationalizes fairness and robustness)}
The adversarial probe constructs ten controlled profile pairs that differ only in a single protected attribute (name, summary, or metadata) and reports the proportion of pairs for which the top-ranked item changes, the maximum score shift, and the disparate-impact ratio (\dir{}) on the experience and remote-policy dimensions.
The probe is narrow by design; it is an engineering probe, not a demographic-fairness audit.
The user-facing interpretation is the system's stability under small, semantically irrelevant edits: a user who makes a typo in a summary should not see their ranking change.

\leadpara{User study stance}
We do not report a live user study in this paper.
The three-week submission window, the small corpus, and the absence of an {IRB}-approved participant pool make a powered user study infeasible before the deadline; we discuss in \secref{sec:8} the study we plan to run after submission.
The evaluation reported here is the strongest evaluation that the prototype, the corpus, and the schedule together support: an end-to-end probe of each design principle, with metrics that map to a user-facing property, on a controlled corpus, with all artifacts released.
